
\medskip
\phantomsection\label{ex:sga3-VIB-5.8.4}%
\noindent\textbf{Example 5.8.4.---}
Retain the notation of the preceding example. Using the
description of \(C\) as the set of sequences
\[
  (u_n)_{n\in\mathbb N},
  \qquad u_n\in\mathbb Z/2\mathbb Z,
\]
one sees that \(C\) with one point \(p\) removed is homeomorphic
to a countable disjoint union of copies of \(C\). Indeed, using
the group structure of \(C\), for example, translation reduces us
to the case in which \(p=0\), the zero sequence. Then
\[
  C-\{0\}=\coprod_{i\in\mathbb N}C_i,
\]
where
\[
  C_i=
  \bigl\{(u_n)_{n\in\mathbb N}\mathrel{\big|}
    u_i=1\ \text{and}\ u_j=0\ \text{for}\ j<i\bigr\},
\]
and every \(C_i\) is homeomorphic to \(C\). It follows by
induction on \(\lvert F\rvert\) that, for every nonempty finite
subset \(F\) of \(C\), the space \(C-F\) is homeomorphic to a
countable disjoint union of copies of \(C\), and hence to \(C\)
with one point removed.

For each \(F\subset C\) of cardinality \(2\), let \(C(F)\) be
another copy of \(C\), denote by \(q_F\) the point \(0\) of
\(C(F)\), and choose a homeomorphism
\[
  \phi_F:C(F)-\{q_F\}\xrightarrow{\sim}C-F.
\]
Let \(X'\) be obtained by gluing every \(C(F)\) to \(C\) by means
of \(\phi_F\). Then \(X'\) is a nonseparated locally Boolean
\(k\)-scheme. Moreover, the construction shows that every locally
constant function
\[
  f:X'\longrightarrow k
\]
is constant. Indeed, if \(x,y\in C\) and \(F=\{x,y\}\), then every
neighborhood of \(q_F\) meets every neighborhood of \(x\) and of
\(y\). Thus
\[
  f(x)=f(q_F)=f(y).
\]
If \(F'=\{z,t\}\), with \(z,t\in C\), then likewise
\[
  f(q_{F'})=f(z)=f(q_{\{z,x\}})=f(x).
\]
Consequently, \(X'\) is connected.

Every point \(x\in X'\) also has a pointed neighborhood
homeomorphic to \((C,0)\). More precisely, for every \(F\), fix a
homeomorphism of pointed topological spaces
\[
  h_F:(C,0)\xrightarrow{\sim}(C(F),q_F),
\]
and let \(T\) be the group of translations of \(C\). If
\(x=q_F\), use \(h_F\); if \(x\in C\), the translation
\(t_x\in T\) is a homeomorphism
\[
  t_x:(C,0)\xrightarrow{\sim}(C,x),
\]
and it is the unique element of \(T\) with this property.%
\footnote[1]{\label{note:sga3-VIB-5.8.4-Xprime-hF}Translator's
note: in these two passages the French source prints \(X\) in
place of the \(X'\) just constructed, and \(q_F\) in place of the
homeomorphism \(h_F\) just defined. The English uses \(X'\) and
\(h_F\), respectively.}

Let \(L\) be the free group generated by the \(h_F\), and set
\[
  G=T*L,
\]
the free product (coproduct) of \(T\) and \(L\). For
\[
  h\in\{h_F\}_{\lvert F\rvert=2}\cup T,
\]
let \(\sigma(h)\) denote the corresponding generator of \(G\),
and let \(S(h)\) and \(B(h)\) be respectively the source and the
target of \(h\). It is convenient also to set
\[
  \sigma(h_F^{-1})=\sigma(h_F)^{-1},\qquad
  S(h_F^{-1})=B(h_F),\qquad
  B(h_F^{-1})=S(h_F),
\]
and to let \(E\) be the set consisting of the elements of \(T\)
and all \(h_F,h_F^{-1}\).

On the product space
\[
  P=G\times X',
\]
where \(G\) carries the discrete topology, consider the
equivalence relation generated by
\[
  \bigl(g\sigma(h),x\bigr)\sim\bigl(g,h(x)\bigr)
  \qquad\text{when }x\in S(h),
\]
for every \(h\in E\), and let \(Z\) be the quotient space. Thus
\(Z\) is obtained by gluing open subsets of the disjoint union
\[
  \coprod_{g\in G}\{g\}\times X'.
\]
Consequently, the saturation \(\widetilde\Omega\) of every open
subset \(\Omega\) of \(P\) is open (cf.\ [BTop], I, \S5.1,
Example~2). Explicitly, since every open subset of \(P\) is the
union of its intersections with the slices
\(\{g\}\times X'\), it is enough to consider
\(\{1\}\times W\), with \(W\) open in \(X'\). Its saturation is
the union of \(\{1\}\times W\) and the sets
\[
  \{\sigma(h_1)\}\times
  h_1^{-1}\bigl(W\cap B(h_1)\bigr),
\]
\[
  \{\sigma(h_1)\sigma(h_2)\}\times
  h_2^{-1}\Bigl(
    h_1^{-1}\bigl(W\cap B(h_1)\bigr)\cap B(h_2)
  \Bigr),
\]
and so on, for all finite sequences
\(h_1,\ldots,h_n\) of elements of \(E\). It is therefore open,
and the projection
\[
  \pi:P\longrightarrow Z
\]
is open.

Moreover, the word
\[
  \sigma(h_1)\cdots\sigma(h_n)
\]
is reduced in \(G\), unless one of the \(h_i\) is the identity of
\(T\), two consecutive \(h_i\)'s belong to \(T\), or one of the
pairs \((h_i,h_{i+1})\) is
\[
  (h_F,h_F^{-1})
  \quad\text{or}\quad
  (h_F^{-1},h_F).
\]
Suppose, therefore, that \(x\in X'\) and that
\[
  \beta=
  \bigl(
    \sigma(h_1)\cdots\sigma(h_n),
    h_n^{-1}\cdots h_1^{-1}(x)
  \bigr)
\]
belongs to \(\{1\}\times X'\). We may then suppose that every
\(h_i\) is a translation \(t_i\). In that case,
\[
  \sigma(t_1\cdots t_n)=1
\]
implies \(t_1\cdots t_n=1\), and hence
\(\beta=(1,x)\). The equivalence relation is compatible with the
action of \(G\) on \(P=G\times X'\) by left translation on the
first factor. It follows that, for every \(g\in G\), the
restriction of \(\pi\) to \(\{g\}\times X'\) is injective.

Now let \(z\in Z\) and choose a representative \((g,x)\in P\).
By the preceding discussion,
\[
  U=\pi\bigl(\{g\}\times X'\bigr)
\]
is an open neighborhood of \(z\), and the continuous map
\[
  \{g\}\times X'\longrightarrow U
\]
induced by \(\pi\) is open and bijective, hence a homeomorphism.
Thus \(Z\) is locally isomorphic to \(X'\), and hence also to
\(C\); it is therefore again a locally Boolean \(k\)-scheme.

Finally, \(G\) acts transitively on \(Z\). Indeed, every \(z\in Z\)
is \(G\)-conjugate to an element of the form
\(\pi(1,x)\), so it is enough to show that every
\((1,x)\in P\) is equivalent to some \((\sigma(h),0)\). If
\(x\in C\), take \(h=t_x\); if \(x=q_F\), take \(h=h_F\). Thus
\(Z\) is a locally Boolean \(k\)-scheme endowed with a transitive
action of the discrete group \(G\), but \(Z\) is not separated.
